\chapter{Kompaktlık}

Kompaktlık, sonsuz yapılarda ``sınırlılık'' fikrinin topolojik soyutlamasıdır.
Açık örtü tanımı temel alınarak incelenmekte; varyantlar ve lokal kompaktlık da ele alınmaktadır.

%-------------------------------------------------------------------
\section{Konu}

\subsection{Sezgisel Giriş — ``Gizli Sonluluk''}

Kompaktlığı tek cümlede özetlemek gerekirse: \emph{sonsuz bir uzayın, sonlu bir
uzay gibi davranmasıdır.} Sonlu bir kümede her şey kolaydır; kompaktlık bu
``sonluluk konforunu'' sonsuz uzaylara taşıyan köprüdür: uzayı ne kadar çok açık
kümeyle örtersek örtelim, \textbf{her zaman sonlu sayıda tanesi yeter.}

\begin{figure}[ht]
  \centering
  \begin{tikzpicture}[font=\small,
    kapsul/.style={rounded corners=6pt, thick, fill opacity=0.18, draw opacity=0.9},
    etiket/.style={font=\scriptsize}]
  % --- Ust: sonsuz acik ortu ---
  \node[etiket, align=left] at (-1.7,1.05) {Açık örtü\\(sonsuz olabilir)};
  \draw[very thick] (0,0.6) -- (8,0.6);
  \node[left] at (0,0.6) {$X$};
  \foreach \x/\c in {0.2/blue, 1.4/red, 2.6/teal, 3.8/orange, 5.0/violet, 6.2/blue, 7.0/red} {
    \fill[\c!60, fill opacity=0.18, rounded corners=5pt] (\x,0.35) rectangle ++(1.6,0.5);
    \draw[\c!70!black, thick, rounded corners=5pt] (\x,0.35) rectangle ++(1.6,0.5);
  }

  % --- ok ---
  \draw[-{Stealth[length=3mm]}, very thick] (4,-0.05) -- (4,-0.7)
    node[midway, right=2pt, etiket] {kompakt $\Rightarrow$ sonlu alt-örtü};

  % --- Alt: sonlu alt-ortu ---
  \node[etiket, align=left] at (-1.7,-1.45) {Sonlu alt-örtü\\(3 küme yeter)};
  \draw[very thick] (0,-1.4) -- (8,-1.4);
  \node[left] at (0,-1.4) {$X$};
  \foreach \x/\c in {0.2/blue, 2.8/teal, 5.6/violet} {
    \fill[\c!60, fill opacity=0.22, rounded corners=5pt] (\x,-1.65) rectangle ++(2.6,0.5);
    \draw[\c!70!black, very thick, rounded corners=5pt] (\x,-1.65) rectangle ++(2.6,0.5);
  }
\end{tikzpicture}

  \caption{Açık örtü $\to$ sonlu alt-örtü: uzayı sonsuz çoklukta açık kümeyle
    örtsek bile, kompaktsa sonlu sayıda tanesi onu kapatmaya yeter.}
\end{figure}

\begin{sezgi}
``Açık örtü'' uzayı kaplayan açık küme ailesidir; ``sonlu alt-örtü'' bu aileden
seçilmiş, hâlâ kaplamayı başaran sonlu bir parçadır. Kompakt olmak, bunu
\emph{her} örtü için yapabilmektir. Bu özellik, analizdeki maksimum değer
teoremi, düzgün süreklilik ve Heine-Borel'in arkasındaki sessiz kahramandır.
\end{sezgi}

\subsection{Kompaktlık Tanımı}

\begin{definition}[Kompakt Uzay]
$(X, \tau)$ topolojik uzayı \textbf{kompakt} ise:
Her açık örtü $\{U_\alpha\}$ için sonlu alt-örtü vardır:
\[
X \subseteq \bigcup_{\alpha} U_\alpha
\implies
\exists \text{ sonlu } I:\; X \subseteq \bigcup_{\alpha \in I} U_\alpha.
\]
\end{definition}

\subsection{Neden ``kapalı ve sınırlı'' yetmez? — Karşı-örnekler}

Reel doğruda kompaktlık tam olarak ``kapalı + sınırlı''ya denktir (Heine-Borel).
Bu denkliğin neden \emph{hem} kapalılığa \emph{hem} sınırlılığa ihtiyaç
duyduğunu en iyi açık aralık gösterir.

\begin{figure}[ht]
  \centering
  \begin{tikzpicture}[font=\small, etiket/.style={font=\scriptsize}]
  % open interval (0,1) on a number line scaled to width 9
  \def\W{9}
  \draw[very thick] (0,0) -- (\W,0);
  \draw[fill=white, thick] (0,0) circle (2.2pt);   % open endpoint 0
  \draw[fill=white, thick] (\W,0) circle (2.2pt);  % open endpoint 1
  \node[below=3pt] at (0,0) {$0$};
  \node[below=3pt] at (\W,0) {$1$};
  \node[etiket, align=center] at (\W/2,3.15)
    {Açık örtü $U_n=\left(\tfrac{1}{n},\,1\right)$, \; $n=2,3,4,\dots$};

  % nested intervals U_n = (1/n, 1): left endpoint at W/n
  \foreach \n/\y/\c in {2/0.55/blue, 3/0.95/teal, 4/1.35/orange, 5/1.75/violet, 6/2.15/red} {
    \pgfmathsetmacro\lx{\W/\n}
    \draw[\c!75!black, very thick] (\lx,\y) -- (\W,\y);
    \draw[\c!75!black, fill=white, thick] (\lx,\y) circle (1.8pt);
    \node[\c!60!black, etiket, left=1pt] at (\lx,\y) {$U_{\n}$};
  }
  % arrow showing escape toward 0
  \draw[-{Stealth[length=2.6mm]}, red!70!black, thick]
    (2.0,2.45) -- (0.35,2.45);
  \node[red!70!black, etiket, align=center] at (4.6,2.5)
    {sol uç $\tfrac{1}{n}\to 0$ ama $0$'a hiç ulaşmaz};

  \node[etiket, align=center, fill=red!6, rounded corners, inner sep=4pt]
    at (\W/2,-1.0)
    {Hiçbir \emph{sonlu} alt-örtü $0$'a yaklaşan kısmı kapatamaz
     $\;\Rightarrow\;(0,1)$ kompakt \textbf{değil}.\\[2pt]
     Kapalı $[0,1]$'de $0$ dahil olduğundan sonlu alt-örtü bulunur $\Rightarrow$ kompakt.};
\end{tikzpicture}

  \caption{Kaçan örtü: $(0,1)$ açık aralığını $U_n=(1/n,1)$ ailesi örter; sol
    uçlar $0$'a yaklaşır ama $0$'a hiç ulaşmaz, bu yüzden hiçbir sonlu alt-örtü
    işe yaramaz.}
\end{figure}

\begin{karsiornek}
\textbf{$(0,1)$ kompakt değil.} $U_n = \left(\tfrac{1}{n}, 1\right)$ ailesi
$(0,1)$'i örter. Sonlu bir alt-aile $\{U_{n_1}, \dots, U_{n_k}\}$ seçersek, en
büyük indis $N=\max n_i$ için birleşim yalnızca $\left(\tfrac{1}{N}, 1\right)$
olur; $0$'a yakın noktalar açıkta kalır. Sonlu alt-örtü yoktur. Kapalı $[0,1]$'de
$0$ dahil olduğundan sonlu alt-örtü her zaman bulunur.

\medskip
\textbf{$\mathbb{R}$ kompakt değil.} $V_n = (-n, n)$ ailesi $\mathbb{R}$'yi örter
ama hiçbir sonlu alt-aile sınırsız $\mathbb{R}$'yi kaplayamaz. Burada sorun
\emph{sınırsızlık}.
\end{karsiornek}

\subsection{Kompaktlık Varyantları}

\begin{table}[h]
\centering
\begin{tabular}{ll}
\toprule
\textbf{Varyant} & \textbf{Tanım} \\
\midrule
Kompakt & Her açık örtünün sonlu alt-örtüsü \\
Sayılabilir kompakt & Her sayılabilir örtünün sonlu alt-örtüsü \\
Ardışımlı kompakt & Her dizi yakınsak alt-dizi içerir \\
Sözde kompakt & Her sürekli $f: X \to \mathbb{R}$ sınırlıdır \\
Lindelöf & Her açık örtünün sayılabilir alt-örtüsü \\
$\sigma$-kompakt & Sayılabilir kompakt kümelerin birleşimi \\
Lokal kompakt & Her noktanın kompakt komşuluğu var \\
\bottomrule
\end{tabular}
\caption{Kompaktlık varyantları}
\end{table}

\begin{figure}[ht]
  \centering
  \begin{tikzpicture}[font=\small, node distance=12mm and 18mm,
    kut/.style={draw, rounded corners, fill=blue!6, inner sep=6pt, align=center},
    vurgu/.style={draw, rounded corners, fill=green!8, inner sep=6pt, align=center, very thick}]
  \node[vurgu] (k) {Kompakt};
  \node[kut, above right=of k] (lind) {Lindelöf};
  \node[kut, below right=of k] (say) {Sayılabilir\\kompakt};
  \node[kut, right=28mm of k] (ard) {Ardışımlı\\kompakt};

  \draw[-{Stealth[length=2.6mm]}, thick] (k) -- (lind);
  \draw[-{Stealth[length=2.6mm]}, thick] (k) -- (say);
  \draw[-{Stealth[length=2.6mm]}, thick] (ard) -- (say);

  % metric equivalence
  \draw[{Stealth[length=2.4mm]}-{Stealth[length=2.4mm]}, dashed, teal!70!black]
    (k.east) to[bend left=10] node[above, font=\scriptsize, align=center] {metrikte\\ $\Leftrightarrow$} (ard.west);

  \node[font=\scriptsize, align=left, anchor=north west] at (-2.2,-2.0)
    {Oklar: $\Rightarrow$ (her zaman geçer).\\
     Kesik çift ok: yalnız \emph{metrik} uzaylarda denklik.};
\end{tikzpicture}

  \caption{Kompaktlık varyantları arasındaki implikasyonlar: Kompakt
    $\Rightarrow$ Lindelöf ve Kompakt $\Rightarrow$ Sayılabilir kompakt her zaman
    geçer; Kompakt $\Leftrightarrow$ Ardışımlı kompakt yalnız metrik uzaylarda.}
\end{figure}

\textbf{Her zaman geçen implikasyonlar:} Kompakt $\Rightarrow$ Sayılabilir
kompakt, Kompakt $\Rightarrow$ Lindelöf, Ardışımlı kompakt $\Rightarrow$
Sayılabilir kompakt. \textbf{Metrik uzaylarda} üç kavram (kompakt, sayılabilir
kompakt, ardışımlı kompakt) çakışır.

%-------------------------------------------------------------------
\section{Teoremler}

\begin{theorem}[Heine-Borel]
$\mathbb{R}^n$'de kompakt $\Leftrightarrow$ kapalı ve sınırlı.
\end{theorem}

\begin{proof}[İspat eskizi]
($\Leftarrow$) $[a,b]$ aralığını ele alalım ve sonlu alt-örtüsü olmayan bir açık
örtü bulunduğunu varsayalım. Aralığı ikiye böleriz; en az bir yarısının da sonlu
alt-örtüsü yoktur. Bunu yineleyerek uzunluğu $0$'a giden iç içe kapalı aralıklar
dizisi elde ederiz; kesişimleri tek bir $p$ noktasıdır. $p$'yi içeren bir örtü
elemanı $U$, yeterince küçük aralıkları tümüyle içerir --- çelişki. $\mathbb{R}^n$
için Tychonoff ile $[a,b]^n$'e geçilir. ($\Rightarrow$) Kompakt küme Hausdorff
uzayda kapalıdır; sınırlılık, $\{(-n,n)^n\}$ örtüsünün sonlu alt-örtüsünden gelir.
\end{proof}

\begin{theorem}
Kompakt $+$ Hausdorff $\Rightarrow$ Normal (T4).
\end{theorem}

\begin{proof}[İspat eskizi]
$A, B$ ayrık kapalı kümeler olsun (kompakt uzayda kapalı $\Rightarrow$ kompakt).
Sabit $a\in A$ için Hausdorff'luk her $b\in B$'yi $a$'dan açık $U_b \ni a$,
$V_b \ni b$ ile ayırır. $\{V_b\}$, kompakt $B$'yi örter; sonlu alt-örtü
$V_{b_1},\dots,V_{b_k}$ alıp $U_a=\bigcap U_{b_i}$, $V_a=\bigcup V_{b_i}$ koyarız.
Burada \emph{sonlu} kesişimin açık kalması kritiktir --- bu yüzden $B$'nin kompakt
olması şarttır. Sonra $\{U_a\}_{a\in A}$ kompakt $A$'yı örter; yine sonlu
alt-örtüden $A$ ile $B$'yi ayıran açıklar kurulur.
\end{proof}

\begin{theorem}
Kompakt uzaydan Hausdorff uzaya sürekli bijeksiyon $\Rightarrow$ homeomorfizma.
\end{theorem}

\begin{theorem}[Tychonoff]
Keyfi kompakt uzayların kartezyen çarpımı kompakttır.
\end{theorem}

\begin{theorem}[Alexandroff Tek-Nokta Kompaktifikasyonu]
Her lokal kompakt Hausdorff uzay $X$ için $X^* = X \cup \{\infty\}$ kompakt
Hausdorff'tur. (Örn.\ $\mathbb{R}^* \cong S^1$, $(\mathbb{R}^2)^* \cong S^2$.)
\end{theorem}

\begin{theorem}[Kompakt görüntü]
$f: X \to Y$ sürekli ve $X$ kompakt ise $f(X)$ kompakttır.
\end{theorem}

\begin{proof}[İspat eskizi]
$\{V_\alpha\}$, $f(X)$'i örten açık aile olsun. Süreklilikten
$\{f^{-1}(V_\alpha)\}$, $X$'i örten açık ailedir; $X$ kompakt olduğundan sonlu
alt-örtü $f^{-1}(V_{\alpha_1}),\dots,f^{-1}(V_{\alpha_k})$ vardır. O hâlde
$V_{\alpha_1},\dots,V_{\alpha_k}$, $f(X)$'i örter.
\end{proof}

%-------------------------------------------------------------------
\section{Algoritmalar}

\subsection{Sonlu Uzayda Kompaktlık}

Sonlu uzay her zaman kompakttır: her örtü zaten sonlu.
Algoritma: $O(1)$.

\subsection{analyze\_compactness\_variants — Tag Tabanlı Çıkarım}

\begin{algorithm}
\caption{AnalyzeVariants$(X, \tau)$}
\begin{algorithmic}[1]
\If{$X$ sonlu}
    \State tüm varyantlar $\leftarrow$ True; \Return VariantProfile(...)
\EndIf
\If{\texttt{'compact'} $\in$ tags}
    \State countably\_compact, sequentially\_compact, lindelof $\leftarrow$ True
\EndIf
\If{\texttt{'metric'} $\wedge$ \texttt{'compact'} $\in$ tags}
    \State sequentially\_compact $\leftarrow$ True
\EndIf
\end{algorithmic}
\end{algorithm}

%-------------------------------------------------------------------
\section{pytop API}

\begin{lstlisting}[language=Python]
from pytop import (
    is_compact,
    is_lindelof,
    is_locally_compact,
    naturals_cofinite,
    one_point_compactification_of_reals,
    one_point_compactification_of_N,
)
from pytop.compactness_variants import (
    is_countably_compact,
    is_sequentially_compact,
    analyze_compactness_variants,
)
\end{lstlisting}

\texttt{analyze\_compactness\_variants(space)} $\to$ \texttt{Result}; \texttt{.value} bir \texttt{dict} içerir.

%-------------------------------------------------------------------
\section{Örnekler}

\subsection*{Örnek 5.1 — Sonlu Uzaylar: Otomatik Kompakt}

\begin{lstlisting}[language=Python]
from pytop import sierpinski_space, discrete_topology, finite_chain_space, is_compact

s = sierpinski_space()
print("Sierpinski compact?", is_compact(s).status)
d = discrete_topology(1, 2, 3)
print("Discrete(3) compact?", is_compact(d).status)
c = finite_chain_space(3)
print("Chain(3) compact?", is_compact(c).status)
\end{lstlisting}

\begin{verbatim}
Sierpinski compact?      true
Chain(3) compact?        true
Discrete(3) compact?     true
\end{verbatim}

\subsection*{Örnek 5.2 — Gerçek Doğru $\mathbb{R}$: Kompakt Değil}

\begin{lstlisting}[language=Python]
rl = real_line_metric()
print("compact?", is_compact(rl).status)
print("lindelof?", is_lindelof(rl).status)
print("locally_compact?", is_locally_compact(rl).status)
\end{lstlisting}

\begin{verbatim}
compact?          false
lindelof?         true
locally_compact?  conditional
\end{verbatim}

\subsection*{Örnek 5.3 — analyze\_compactness\_variants: Sierpiński}

\begin{lstlisting}[language=Python]
r = analyze_compactness_variants(s)
for key, val in r.value.items():
    if hasattr(val, 'status'):
        print(f"  {key:<25}: {val.status}")
\end{lstlisting}

\begin{verbatim}
  representation           : finite
  countably_compact        : true
  sequentially_compact     : true
  pseudocompact            : true
  lindelof                 : true
\end{verbatim}

\subsection*{Örnek 5.4 — Sonsuz Ama Kompakt: Kosonlu Topoloji}

Kompaktlık sınırlılığa değil, \emph{örtü davranışına} bağlıdır. Sonsuz bir küme
üzerindeki kosonlu topoloji buna örnektir: bir açık örtüde tek bir eleman
seçtiğimizde dışarıda yalnız sonlu çoklukta nokta kalır; onları da sonlu sayıda
elemanla kapatırız.

\begin{lstlisting}[language=Python]
from pytop import naturals_cofinite, is_compact, is_lindelof
from pytop.compactness_variants import is_countably_compact

nc = naturals_cofinite()
print("compact?", is_compact(nc).status)
print("lindelof?", is_lindelof(nc).status)
print("countably?", is_countably_compact(nc).status)
\end{lstlisting}

\begin{verbatim}
compact?   true
lindelof?  true
countably? true
\end{verbatim}

Sonsuz olmasına karşın kompakt --- ``kompakt $=$ sonlu'' sezgisinin neden yalnız
\emph{metrik} uzaylarda işlediğini gösterir.

\subsection*{Örnek 5.5 — Tek-Nokta Kompaktifikasyon (Alexandroff)}

Lokal kompakt ama kompakt olmayan bir uzaya tek bir $\infty$ noktası ekleyerek onu
kompakt hâle getirebiliriz. $\mathbb{R}$ için sonuç çembere ($S^1$) denktir.

\begin{lstlisting}[language=Python]
from pytop import (
    one_point_compactification_of_reals,
    one_point_compactification_of_N,
    is_compact,
)

r_star = one_point_compactification_of_reals()   # R u {inf} ~= S^1
n_star = one_point_compactification_of_N()
print("R* compact?", is_compact(r_star).status)
print("N* compact?", is_compact(n_star).status)
\end{lstlisting}

\begin{verbatim}
R* compact? true
N* compact? true
\end{verbatim}

$\mathbb{R}$ kompakt değildi (Örnek 5.2); tek bir nokta eklemek onu kompakt yaptı.

%-------------------------------------------------------------------
\section{Alıştırmalar}

\subsection*{Kodlama}

\begin{enumerate}[label=K\arabic*.]
\item \texttt{finite\_chain\_space(5)} ve \texttt{discrete\_topology(1,2,3,4,5)} için
      \texttt{is\_compact} karşılaştırın.
\item \texttt{naturals\_cofinite()} için \texttt{is\_compact}, \texttt{is\_lindelof},
      \texttt{is\_countably\_compact} hesaplayın.
\item \texttt{closed\_unit\_interval\_metric()} ile \texttt{real\_line\_metric()} için
      \texttt{analyze\_compactness\_variants} çalıştırıp karşılaştırın.
\item \texttt{real\_line\_metric()} kompakt değildi.
      \texttt{one\_point\_compactification\_of\_reals()} çıktısının
      \texttt{is\_compact} durumunu kontrol edip tek noktanın kompaktlığı nasıl
      kurtardığını açıklayın.
\end{enumerate}

\subsection*{Teori}

\begin{enumerate}[label=T\arabic*.]
\item Kompakt + sürekli $\Rightarrow$ görüntü kompakt olduğunu ispatlayın.
\item Heine-Borel teoreminin neden geçerli olduğunu: $[0,1]$ neden kompakt,
      $(0,1)$ neden kompakt değil? (İpucu: kaçan örtü $U_n=(1/n,1)$.)
\item Kompakt $+$ Hausdorff $\Rightarrow$ Normal ispatında, $B$'yi örten
      $\{V_b\}$ ailesinden sonlu alt-örtü çekme adımının neden kompaktlığa
      ihtiyaç duyduğunu açıklayın.
\end{enumerate}
